\section{Introduction}

Single-cell RNA sequencing has grown over a decade from dozens of cells to atlases containing millions. The operation that makes this growth useful is \emph{batch integration}: the alignment of data across laboratories, sequencing platforms, dissociation protocols, and collection time points so that downstream analysis reflects biology rather than technical variation. Every integrator that the community has adopted --- Harmony, Seurat, scVI, scANVI, Scanorama, fastMNN, BBKNN, LIGER, scGen, scDML, scDREAMER, SCIDRL, scCRAFT, STACAS, CellANOVA --- occupies a point on a single tradeoff: the more aggressively batch variation is removed, the greater the risk that genuine biological variation is removed with it. This tradeoff is called \emph{overcorrection}, and it is the central tension in the integration literature \citep{scDML2023,scDREAMER2023,SCIDRL2022,scCRAFT2025,STACAS2024,CellANOVA2024,RBET2025,scIB-E2025}.

Rare cell populations are the first casualty of overcorrection. Their transcriptomic signal is weak relative to the scale of typical batch effects; their cells occupy a small volume in the embedding; and when integration aligns batches, they can be displaced into adjacent abundant populations and cease to exist as distinct biological entities. Multiple high-impact studies have documented this phenomenon on specific previously-discovered rare populations: pancreatic epsilon cells, pulmonary ionocytes, Schwann cells, activated stellate cells, OFFx bipolar cells in macaque retina, TCF7$^+$ stem-like CD8, LAMP3$^+$ mature dendritic cells, FOLR2$^+$ tumor-associated macrophages. A 2026 \emph{Nature Computational Science} review summarised that despite a decade of effort the field still oscillates between overcorrection and undercorrection. And yet the problem is not that overcorrection goes undetected --- it is that it is detected \emph{only} for populations the analyst already knew about.

The asymmetry has a second face: for five years the field has attempted to solve this problem and has not.  The failure is not for lack of attention.  It is structural.  Detecting loss of an unannotated rare subpopulation requires an evaluation framework; building such a framework requires a definition of success, which in turn requires ground truth for what was supposed to be preserved; and acquiring ground truth for populations that are by construction unannotated is precisely the step the framework was meant to replace.  No ground truth begets no evaluation, which begets no failure signal, which begets no motivation to develop a protective algorithm.  The loop is closed.  It is not, however, unbreakable: normal batch correction produces convergence (same-type cells merged across batches, cluster still present) whereas erasure of a rare subpopulation produces dispersion (cluster members scattered into unrelated large clusters), and these two patterns are structurally different in the data itself.  The signal exists; what was missing was a framework to extract it without external labels.

% Measured variant of the page-1 framing paragraph.
% For task t-mozveh1b (Mathematics) and t-mozv8oul (Biological Science).
% Same content, lower-temperature prose — pairs with the aggressive variant so
% editors can pick one. References Math Theorem 2.1 and the exchangeability bound (Lemma S1).

\paragraph*{The limits of label-based evaluation.}
Standard evaluation of single-cell integration relies on quantities---$\mathrm{ARI}$, $\mathrm{NMI}$, $\mathrm{ASW}$, and related scores in scIB, scIB-E, and RBET---that are functionals of annotator labels. A formal consequence, stated as Theorem~\ref{thm:labelblind}, is that every label-based functional is invariant under a class of local coupling deformations that can eliminate an unannotated rare component by folding it onto an adjacent annotated population. In other words, the failure mode the field has most worried about for at least five years---loss of rare subpopulations that no atlas has yet named---sits outside the measurable sets of its own evaluation apparatus. The failure is not that existing metrics are imprecise; it is that they cannot, in principle, observe the events of interest. This matters downstream: the inference ``integrated data are safe for discovery because no label-based metric dropped'' conflates \emph{absence of evidence} with \emph{evidence of absence} along a dimension the evaluation framework cannot read. A productive response is not to reject integration but to evaluate it with instruments that do not presuppose the annotation. Accordingly, we propose REAL (Rare-subpopulation Erasure under Alignment, Label-free), which assesses integration through a label-free, geometrically-defined notion of local-mass preservation and is calibrated by a held-out known-rare protocol under an explicit exchangeability bound (Lemma~\ref{lem:exchangeability}). REAL does not settle questions of biological novelty; it marks where the evaluation framework has been blind.


% Introduction paragraph 3 (≈250w) — Meno-problem framing.
% For task t-mozv8oul (Biological Science) and BioSci's locked abstract/Intro-P3 slot.
% Anchors: Math Assumption (A4) / Lemma S1-exchangeability (bound), framework name REAL,
% motif set W, protection RareShield, per-seed score LRS(w).

\paragraph*{A Meno problem in empirical science.}
Plato's \emph{Meno} asks how one can search for what one does not know: if the object of inquiry is unknown, it cannot be recognised even if encountered. The preservation of previously unannotated rare cell subpopulations during single-cell batch integration instantiates this paradox in a concrete form. Every existing evaluation---scIB, scIB-E, RBET, CellANOVA---certifies preservation against an annotation. If a subpopulation has never been annotated, no score moves when it is destroyed, and no score moves when it is preserved; the object is simply not in the evaluator's vocabulary. Our resolution is to replace the \emph{semantic} criterion (does the population have a name?) with a \emph{syntactic} criterion (is there reproducible structure across independent observations?). Reproducibility is a label-free property, witnessed geometrically by local-density motifs $w \in \mathcal{W}$ that recur across batches and persist under admissible batch-effect deformations, and quantified by a per-motif survival score $\mathrm{LRS}(w)$. A motif whose cross-batch reproducibility predates integration but not its integrated image is diagnostic of erasure, regardless of whether any name has been attached to it. We validate the criterion on held-out known rare types (pancreatic $\epsilon$, pulmonary ionocytes, LAMP3$^+$ mregDC, TPEX, FOLR2$^+$ macrophages) by hiding their labels, confirming detection, and transferring the estimate to genuinely unannotated subpopulations under an exchangeability bound (Lemma~\ref{lem:exchangeability}). REAL detects; RareShield protects. The bound, not the field's habits, licenses the transfer.


% ---------------------------------------------------------------------
% Immunology paragraph for the Introduction of the CNS submission.
% ~500 words. Contributed by Immunology agent (agent-mozus3vv) for
% task t-mozv8ou4 (follow-up: Biological Science asked for this
% explicitly as workspace/handoffs/immuno_intro_para.tex).
% Target: sections/01_introduction.tex, probably after the general
% motivation paragraph on integration + overcorrection and before
% the paragraph that introduces our framework (LRS / RareShield).
% ---------------------------------------------------------------------

The stakes of this failure mode are not uniformly distributed across
the cell-biology landscape; they are most concentrated in the immune
compartment, which is also the compartment where the majority of
integration-atlas effort now accumulates. Three structural properties
of immunology make unannotated rare subpopulations especially
vulnerable to batch-integration overcorrection. First, the immune
system is a lineage continuum: committed lineages pass through
progenitor, activation, and terminal states whose rare intermediates
(pre-DC, transitional B, progenitor-exhausted CD8) sit at the joints
of an explicitly hierarchical tree \citep{villani2017single,
cancro2020abcs}. Because integration methods penalize between-batch
dispersion in a shared latent space, these joint populations --- the
populations that are rare precisely because they are transitional ---
are the populations with the highest expected displacement under
batch alignment. Second, an immune cell's transcriptome is tissue-,
antigen-, and context-dependent to a degree rarely matched elsewhere:
a circulating memory CD8 T cell and its tissue-resident counterpart
share at most 60\% of their differentially expressed genes, and
mucosal-associated invariant T cells (MAIT) in blood, gut, and lung
cluster as effectively distinct populations despite identical TCR
identity \citep{provine2020mait, dominguez2022crosstissue}. Standard
batch covariates (patient, chemistry, sequencing centre) collapse
precisely those axes, rendering biologically informative structure
indistinguishable from technical noise. Third, prevalence itself is
informative: the clinical readout in immunology is often the fold
change in a rare population's abundance across conditions --- MAIT
collapse in severe COVID-19, age-associated B-cell expansion in
systemic lupus erythematosus, CD103$^+$CD39$^+$ tumour-reactive
tissue-resident memory CD8 ratios between checkpoint-blockade
responders and non-responders. Integration methods that normalize
cell-type prevalence across batches to ``correct'' sampling
imbalance obliterate exactly the signal the downstream analysis is
designed to measure \citep{duhen2018coexpression,
caushi2021transcriptional, vandamme2021depletion}.

The historical record amplifies the concern. Nearly every
mechanistic advance in the past decade of immunology began as a
small, ambiguously-identified cluster that predated its nameable
identity. Plasmacytoid dendritic cells, at $<\!0.5\%$ of peripheral
blood and transcriptomically adjacent to B cells, would not have
established the type-I interferon axis of antiviral immunity had
they been absorbed into their abundant neighbour during integration;
the \textit{AXL}$^+$\textit{SIGLEC6}$^+$ ``AS-DC'' population (DC5),
at $\sim\!0.05\%$ of PBMC, rewrote the classification of the
conventional DC lineage only after Villani \etal\ isolated it as a
distinct cluster from several thousand cells
\citep{villani2017single}; tumour-reactive \textit{CD103$^+$CD39$^+$}
tissue-resident memory T cells at $<\!1\%$ of tumour cells now drive
adoptive-cell-therapy patient selection in multiple trials
\citep{duhen2018coexpression}. In every case, the cluster existed
before the biology was understood, and in every case, an integration
pipeline that silently merged the cluster into its abundant
neighbour would have delayed or extinguished the discovery. To the
degree that pan-cancer and pan-tissue atlases now dominate immune
biology \citep{kang2024integrated, cheng2021pancancer,
zheng2021pancancer}, integration is no longer a pre-processing step;
it is the step at which future rare-population discoveries are
either protected or silently lost. A detection framework that
operates without relying on prior annotation, and an integration
strategy that is rare-aware by construction, are therefore not
refinements of an existing toolkit but the prerequisites for the
next generation of immunological discovery.

% ---------------------------------------------------------------------
% Requires: villani2017single, cancro2020abcs, provine2020mait,
% dominguez2022crosstissue, duhen2018coexpression,
% caushi2021transcriptional, vandamme2021depletion, kang2024integrated,
% cheng2021pancancer, zheng2021pancancer. All in
% workspace/handoffs/bib_additions_immunology.bib.
% ---------------------------------------------------------------------


% ============================================================
% Introduction tumor-biology paragraph for sections/01_introduction.tex
% Owner: Tumor Biology Agent (agent-mozurh1r)
% Target: ~350 words, LaTeX-ready, to splice after the "why batch integration matters" paragraph
% and before the detectability-gap statement. Framework name: REAL. Protection: RareShield.
% ============================================================

\paragraph{Why cancer data is the worst case.}
The consequences of silent rare-subpopulation erasure are not uniform across biology. Cancer single-cell data is, by construction, the domain in which the conditions that make a population integration-fragile are stacked most tightly. Three properties co-occur. First, the populations that carry the most therapeutic weight in oncology are structurally rare: drug-tolerant persister cells are fewer than one percent of tumor cells at diagnosis \citep{Rambow2018Cell,Shen2020Cell,Oren2021Nature,Marine2020NatRevCancer}; LAMP3$^+$ mature regulatory dendritic cells are under three percent of tumor dendritic cells \citep{Maier2020Nature}; CCR8$^+$TNFRSF9$^+$ tumor effector regulatory T cells, the target of at least five Phase 1/2 programmes, are typically below one percent of the tumor microenvironment \citep{Plitas2016Immunity,DeSimone2016Immunity,Wang2019Cell,VanDamme2021,Kidani2022PNAS}. Second, these populations sit transcriptomically adjacent to their abundant neighbors: TCF7$^+$ stem-like CD8 T cells emerge from the exhausted CD8 cluster \citep{Miller2019NatImmunol,Siddiqui2019Immunity}; TREM2$^+$ lipid-associated macrophages sit inside the general tumor-associated-macrophage distribution \citep{Molgora2020Cell,Katzenelenbogen2020Cell}; antigen-presenting cancer-associated fibroblasts (apCAF) are routinely mislabeled as immune cells because of HLA-DR expression \citep{Elyada2019CancerDiscov}. Third, cancer cohorts are by construction non-exchangeable across batches: batch in cancer single-cell atlases partially encodes cancer type, treatment history, and disease state. Any integration procedure that enforces isotropy across batches therefore removes biology in proportion to how non-exchangeable the cohorts are.

The compounding effect is direct: the populations whose preservation most matters for precision oncology are the populations whose pre-integration signal is smallest, whose distance to their abundant neighbor is shortest, and whose cross-batch presence is the most asymmetric. Every existing integration benchmark has missed this failure mode because it measures preservation against labels that an earlier annotator already knew, and the populations at issue are frequently populations that annotators had not yet named. Here we resolve the gap: REAL supplies a label-free detectability framework that ranks reproducible rare motifs against a null of non-reproducibility, RareShield supplies a protective integration augmentation that reduces motif loss without compromising major-cell-type integration, and we apply both to the largest pan-cancer immune atlas currently available \citep{Kang2024PanCancerImmune} to estimate how often cancer rare biology has been silently collapsed under standard pipelines.


% Onco-specific Intro paragraph (~140w). For Tumor Biology agent-mozurh1r.
% Complements tumor_bio_intro_para.tex (three-property framing: low abundance + adjacency + non-exchangeable cancer batches).
% Philosophy-level: WHY cancer is the hardest exchangeability case, connecting to Lemma S1.
% Target: sections/01_introduction.tex, to follow tumor_bio_intro_para.tex immediately.

Cancer is the limiting case for the exchangeability bound \Cref{lem:exchangeability}. In non-malignant tissue atlases, batches typically vary along technical axes (donor, chemistry, centre) that are approximately independent of the biology under study, so the bound is tight. Cancer atlases invert that relationship: each batch encodes a disease state --- tumour subtype, treatment history, micro-environmental niche --- as an inseparable covariate of biology itself. Batch-correction isotropy therefore removes genuine biology in direct proportion to the non-exchangeability of batches, and the known-rare-to-unknown-rare transfer that grounds our validation (Lemma~\ref{lem:exchangeability}) is least guaranteed precisely where it is most clinically consequential. This is not a peripheral concern: the populations whose presence is diagnostic of response to immunotherapy or targeted therapy --- TPEX, LAMP3$^{+}$ mregDC, TREM2$^{+}$ LAM, CCR8$^{+}$ eTreg, the p-EMT hybrid, apCAF --- are all born at this junction. A framework that claims pan-cancer-atlas scalability must therefore not only satisfy the minimax envelope of \Cref{thm:minimax} but explicitly disclose the cancer-specific boundary condition under which its exchangeability bound is least tight.


Here we close the gap.  We introduce the \emph{Rare-subpopulation Erasure under Alignment, Label-free} (\textsc{REAL}) framework, which defines a rare subpopulation as a reproducible, multi-batch, high-density region in gene-expression space, and scores any integrator by how well it preserves such regions via a per-seed survival score, the LRS.  REAL does not consult annotation at any stage, which lets it detect loss of populations that were never named.  We derive \textsc{RareShield}, a protective integration procedure that anchors REAL-witnessed motifs in the latent space of any differentiable integrator without requiring labels.  We validate the pair on (i) label-masked positive controls where known rare populations are hidden from the pipeline, (ii) synthetic masked-unknown controls where rare populations are injected and then hidden, (iii) three immunology-anchored experiments (MAIT preservation across tissue batches, eTreg dilution in tumor atlases, pDC loss in severe COVID), (iv) a scale stress test on the Kang 2024 pan-cancer immune atlas with 4.9~million cells, and (v) cross-atlas replication of the flagged motifs on independent tumor atlases.  We close by arguing that the label-based benchmarks of the last decade should be extended with a label-free axis (which we call scIB-U), and that a fraction of the downstream conclusions built on existing integrated atlases should be re-examined under this axis.
